% 01_intro.tex — Introduction (Nature Article lead; Related work folded in).
% Maier voice per ../STYLE.md (plain, precise, short sentences; no poetic register).
% The Introduction is the unnumbered lead right after the abstract (Nature Article
% format), so it carries no \section heading.
% Arc (each point stated ONCE): problem (faithfulness, unmeasurable) -> the three
%  traditions share the blind spot -> the oracle makes it measurable (Fig 1) ->
%  result 1: faithfulness measured -> result 2: the faithfulness escape removed
%  (sampler on), the gap survives -> the semantic gap, both directions, F/S/M ->
%  meaning's reference is behaviour -> contributions and scope.
% Detailed benchmark comparison (BAM/M4/Tracr/InterpBench/MIB), the Atari-environment
% lineage, and the competing definitions of "interpretable" (Barbiero; Marr; Lazebnik)
% live in the Discussion to avoid repeating them here; every \cite key is preserved
% across the two sections (see the whole-paper citation check in the rework commit).
% Numbers: sampler from compare/out/sampler_faithfulness.csv (R7-1); rest from number_audit §1.

Interpretability now matters in practice. Neural networks drive cars, read scans, and
approve loans. Before we trust such a decision, we ask why. Yet on the systems we most
want to explain we cannot check whether the answer is correct. A saliency map for a deep
network can look convincing and can change with the action, and still not be a correct
explanation~\cite{greydanus2018visualizing}. The field of explainable AI has converged on
\emph{faithfulness} as its quality criterion: whether an explanation names a system's
\emph{true} causes~\cite{jacovi2020towards}. The difficulty is that on a real neural
network there is nothing to measure it against. Sanity checks can show that a method fails,
but with no reference map they cannot certify that it is correct~\cite{adebayo2018sanity}.

This blind spot is not specific to deep learning. Jonas and Kording applied standard
neuroscience methods to a microprocessor in which every transistor is
known~\cite{jonas2017could}. Lesions, tuning curves, and connectomics produced rich
structure, but none of it recovered how the chip computes. The same gap recurs across the
three interpretability traditions that rarely cite one another. The neuroscience battery
records and perturbs a system. The attribution tradition asks which inputs mattered for a
decision, through gradient and saliency
maps~\cite{simonyan2014deep,sundararajan2017axiomatic,smilkov2017smoothgrad,springenberg2015striving},
class-activation methods~\cite{selvaraju2017gradcam,chattopadhyay2018gradcampp},
occlusion~\cite{zeiler2014visualizing,fong2017interpretable,petsiuk2018rise}, and surrogate
or game-theoretic attributions~\cite{ribeiro2016why,lundberg2017unified}, applied to Atari
agents to visualise what a policy attends~\cite{greydanus2018visualizing,such2019atarizoo}.
Mechanistic interpretability aims past attribution at the circuit, through path patching and
causal mediation~\cite{vig2020investigating,meng2022locating,wang2022interpretability,conmy2023towards},
alignment search and interchange interventions~\cite{geiger2021causal,geiger2023finding}, and
the sparse dictionaries and probes that read features off the
state~\cite{cunningham2023sparse,alain2017understanding,hewitt2019designing,belrose2023eliciting,nanda2023attribution}.
Across all three, methods are validated against each other or against human intuition rather
than against the mechanism~\cite{atrey2020exploratory}, and even the sparse-dictionary branch
reports that its tuning proxies do not track interpretability~\cite{saebench2025}.
Faithfulness has stood in for understanding, and no one could test the stand-in.

Maier et al. recently built the missing reference. It is a bit-exact, differentiable port
of the Atari 2600 Video Computer System (VCS) that reproduces a real 1977 machine to the
byte and pixel~\cite{maier2025vcs,xitari}. The VCS is software with a known specification,
so the true cause of any output is computable by exact intervention. This \emph{intervention
oracle} (\Sectionname\ref{sec:oracle}) returns the true causal effect for any candidate
cause of any output, measured rather than estimated. It operates over the platform shown in
Fig.~\ref{fig:platform}. On this machine an interpretability claim can be scored right or
wrong. We operationalise interpretability as ground-truth recovery: an explanation is
interpretable to the degree that its claimed structure matches the system's true causal
structure, scored by the oracle.

\begin{figure}[t]
  \centering
  \includegraphics[width=\linewidth]{fig1_platform_oracle.pdf}
  \caption{\textbf{The measurement platform and its ground-truth oracle.} A
  differentiable, bit-exact port of the Atari VCS (left), realized as
  \texttt{jutari} (Julia) $\cdot$ \texttt{jaxtari} (JAX), gives four capabilities no neural
  network grants for free: exact forward, exact interventions $\mathrm{do}(u)$, full state
  observability, and content-path gradients. On it we build the oracle (centre): the primary
  intervention oracle of \Sectionname\ref{sec:oracle}, $\mathrm{do}(u{:=}v')\!\to\!\Delta y(u)$
  by bit-exact re-run, with a content-path gradient $\partial y/\partial u$ as a companion on
  the content path only. Any explanation $\hat{E}$ is scored by the correctness triad (right),
  $\text{right}=\mathrm{F}\wedge\mathrm{S}\wedge\mathrm{M}$, and reported on two axes,
  faithfulness against the oracle versus a plausibility proxy. Faithfulness aggregates
  correlation, sign and magnitude agreement, deletion/insertion area-under-curve, and
  precision at $k$ against the true causal top-$k$ (\Sectionname\ref{sec:triad}). A discrete index --- a
  sprite's position, set by strobe timing --- has a naive gradient that is zero under the hard
  forward map (Prop.~\ref{prop:zero}); for position, index, and event outputs the intervention
  oracle is the \emph{sole} ground truth and gradient methods are scored as methods under test.
  The schematic draws no experiment data. Data and oracle: see Supplement; oracle defined in
  \Sectionname\ref{sec:oracle}.}
  \label{fig:platform}
\end{figure}

We audit $30$ interpretability methods plus the oracle positive control, across $31$
experiments and $257$ per-game records, scoring each on two axes: faithfulness against the
oracle, and a documented plausibility proxy for how convincing it looks. The first result is
the expected one. Causal and intervention methods are faithful; gradient and correlational
saliency methods are not. The contrast is sharpest on discrete sprite-position outputs, where
the naive gradient is provably zero because position is set by strobe timing rather than by a
value it can flow through. There activation patching~\cite{vig2020investigating,meng2022locating}
reaches faithfulness $1.000$ while vanilla saliency~\cite{simonyan2014deep} collapses to
$0.000$ on all six core games~\cite{faithfuldemoP2}, and the leaderboard plots all $31$ methods
against the same ground truth (Fig.~\ref{fig:faithfulness})~\cite{leaderboardP2}. The saliency
method nonetheless carries the higher plausibility score ($0.9$ versus $0.5$): it looks most
convincing exactly where it is least correct. This is the result the field would predict, and
we treat it as an intermediate step, not the main finding.

Faithfulness is not what separates good explanations from bad ones. A reader could dismiss the
zero of the gradient methods as a fixable artifact of a discrete index, so we test the
strongest version of the question and close that objection. We turn the differentiable port's
bilinear sampler on. The sampler interpolates the position read and restores a non-zero
gradient through it, and the gradient methods become faithful too. Vanilla saliency,
grad$\times$input, smoothgrad, and expected gradients reach $0.791$ on Pong, and the same
family reaches $0.529$ on Q*bert (Fig.~\ref{fig:sampler}). The rise is real but not universal:
integrated gradients with a zeros baseline still vanishes on Pong, and a static-sprite game
offers no position to restore. Yet on every game where the gradient becomes faithful, the
method still recovers no meaning, just as no other method does, faithful or not. The restored
gradient concentrates on the position RAM byte that drives the sprite, \texttt{ram[54]} on
Pong. It identifies the correct cause but supplies no meaning for it.

Securing faithfulness is necessary, but it does not solve the main problem. With the same
oracle we ask not only \emph{which} methods are faithful but \emph{what} a faithful method
delivers, and the answer is always the same kind of object: a true causal-effect map or
data-flow graph, never the variable's meaning and never the routine's algorithm. To make this
precise we call an explanation right only if it is \textbf{F}aithful (its claimed causes are
the true ones), \textbf{S}ufficient (it regenerates behavior under held-out interventions),
and \textbf{M}inimal at the algorithmic level
($\mathrm{right} = \mathbf{F}\wedge\mathbf{S}\wedge\mathbf{M}$, \Sectionname\ref{sec:triad}).
The committed data show the gap in both directions (Fig.~\ref{fig:taxonomy}). An automatically
discovered circuit on Breakout is graph-correct and minimal yet reproduces the behavior it
explains only part of the time ($\mathbf{F}=1.0$, $\mathbf{M}=1.0$, $\mathbf{S}=0.44$): a
correct wiring diagram is not the computed function. A learned feature dictionary on Pong
matches a named game variable perfectly (matched fraction $1.0$) while being causally almost
inert ($\mathbf{F}=0.04$): looking like a known variable is not computing it the way the
program does. A concept can be linearly decodable yet provably unused; being present in the
state is not being used. We name this residual the \emph{semantic gap}, and it does not shrink
when faithfulness is made universal.

The reference for meaning lies outside the mechanism. Every semantic label in the study entered
from outside the system, imported from external annotation
sets~\cite{anand2019unsupervised,delfosse2023ocatari} and only \emph{causally verified} on our
substrate, never \emph{discovered} by a method under test. An alignment method that reports
perfect accuracy is therefore right by construction, because it checks a variable we supplied.
Probing shows that information is present, not that it is used~\cite{hewitt2019designing}, and
our oracle is the strictly stronger test. What fixes the meaning is the system's
\emph{behavior}. For software, that behavior is set down in its documentation and data: the
manuals, the specification, and the design. IEEE~610.12 makes these part of the definition of
software itself~\cite{ieee610}. The wiring is in the mechanism, so any faithful method can
recover it. The meaning is not, and no method we tested supplied it.

To our knowledge this is the first cross-tradition audit of interpretability methods on a real,
human-written software artifact whose ground truth is independent of it, scored by an exact
intervention oracle. We claim no unqualified ``first ground-truth calibration,'' since
synthetic known-answer benchmarks and Jonas and Kording's microprocessor
study~\cite{jonas2017could} precede us; the Discussion places our oracle against those
benchmarks and against the competing senses of ``interpretable.'' We read a 1977 chip as a
\emph{necessary-condition screen}, not a predictor of success on learned networks. Here we hold
perfect ground truth and no learning obscures the mechanism, so failure under these conditions
is strong negative evidence that the same method would recover causal mechanisms in harder,
less observable systems. The chip still exhibits the properties that make interpretation hard,
from aliased RAM cells to race-the-beam timing.

We are deliberate about what we do not claim. No method recovered the program's meaning,
algorithm, or design from scratch, and the plausibility axis is a documented proxy, not a
human study. This paper \emph{defines and measures} the semantic gap and locates its reference
in the system's behavior. \emph{Closing} the gap means recovering a system's documentation and
design from the system alone. That is future work, hopefully inspired by this paper. The study
thus offers a reusable ground-truth benchmark for explainable AI, mechanistic interpretability,
AI safety, and software verification, and it reframes the open problem. The hard part of
interpretability is not faithfulness. It is recovering meaning, which remains unsolved even when
faithfulness is achieved. We believe that naming and demonstrating the existence of this gap is
a first step toward closing it.


% BIBTEX-NEEDED manifest (provenance for the SM merge; comments only, no typeset effect).
% Keys cited HERE that may need merging into references.bib: maier2025vcs,
% anand2019unsupervised, hewitt2019designing (plus those already present: jonas2017could,
% atrey2020exploratory, greydanus2018visualizing, xitari, delfosse2023ocatari).
% Keys formerly in this section that now live in the DISCUSSION after the intro rework
% (do NOT re-add them here): barbiero2025neural, bellemare2013arcade, machado2018revisiting,
% dalton2020cule, delfosse2025jaxatari (Atari-environment lineage + "interpretable" variants),
% and the benchmark-comparison keys lindner2023tracr, gupta2024interpbench, yang2019bim,
% m4_2023, mib2025, marr1982vision, lazebnik2002can (08_discussion.tex carries them).
%
% @article{jonas2017could,
%   title   = {Could a Neuroscientist Understand a Microprocessor?},
%   author  = {Jonas, Eric and Kording, Konrad Paul},
%   journal = {PLoS Computational Biology},
%   volume  = {13}, number = {1}, pages = {e1005268}, year = {2017},
%   publisher = {Public Library of Science}, doi = {10.1371/journal.pcbi.1005268}}
%
% @inproceedings{atrey2020exploratory,
%   title     = {Exploratory Not Explanatory: Counterfactual Analysis of Saliency Maps for Deep Reinforcement Learning},
%   author    = {Atrey, Akanksha and Clary, Kaleigh and Jensen, David},
%   booktitle = {International Conference on Learning Representations (ICLR)},
%   year      = {2020}}
%
% @inproceedings{greydanus2018visualizing,
%   title     = {Visualizing and Understanding Atari Agents},
%   author    = {Greydanus, Sam and Koul, Anurag and Dodge, Jonathan and Fern, Alan},
%   booktitle = {International Conference on Machine Learning (ICML)},
%   pages     = {1792--1801}, year = {2018}, organization = {PMLR}}
%
% @inproceedings{anand2019unsupervised,
%   title     = {Unsupervised State Representation Learning in Atari},
%   author    = {Anand, Ankesh and Racah, Evan and Ozair, Sherjil and Bengio, Yoshua and C\^{o}t\'{e}, Marc-Alexandre and Hjelm, R Devon},
%   booktitle = {Advances in Neural Information Processing Systems (NeurIPS)},
%   volume    = {32}, year = {2019}}
%
% @inproceedings{delfosse2023ocatari,
%   title     = {OCAtari: Object-Centric Atari 2600 Reinforcement Learning Environments},
%   author    = {Delfosse, Quentin and Bl{\"u}ml, Jannis and Gregori, Bjarne and Sztwiertnia, Sebastian and Kersting, Kristian},
%   booktitle = {Reinforcement Learning Conference (RLC)}, year = {2024},
%   note      = {arXiv:2306.08649}}
%
% @inproceedings{hewitt2019designing,
%   title     = {Designing and Interpreting Probes with Control Tasks},
%   author    = {Hewitt, John and Liang, Percy},
%   booktitle = {Conference on Empirical Methods in Natural Language Processing (EMNLP)},
%   pages     = {2733--2743}, year = {2019}}
%
% @misc{maier2025vcs,
%   title  = {A Bit-Exact, Differentiable Atari 2600 (Paper 1)},
%   author = {Maier, Andreas and Bayer, Stefan and Krauss, Patrick},
%   year   = {2025}, note = {companion paper; this program}}
